Debt amounts add a second distinction. Note~20 puts US\$253{,}039m of
long-term-debt carrying value under Parent company for 2025, while Note~33's
parent-only balance sheet shows US\$254{,}633m. The US\$1{,}594m difference
equals Note~20's Subsidiaries junior-subordinated-debt subtotal. The same
comparison gives US\$1{,}603m in 2024. This exact two-year arithmetic match
does not establish why a subsidiary-labelled row aligns with the parent-only
balance, or which legal entity owes each instrument. Its accounting and
contractual relation still needs source support.

Timing requires a different measure. Note~20 places US\$14{,}298m of parent
carrying value in its under-one-year bucket; Note~33 reports US\$14.3bn of
parent long-term debt contractually maturing in 2026. They align at the
footnote's displayed precision, but carrying value includes premiums,
discounts, costs and valuation adjustments rather than simply principal
payable. For the separately reported consolidated zero-coupon population,
Note~20 gives US\$18.1bn of carrying value against US\$57.6bn of aggregate
principal at the notes' respective maturities. The firm's US\$435{,}206m
consolidated long-term-debt total and US\$47.6bn of subsidiary debt and
structured notes guaranteed by the parent are also separate from the parent's
own debt balance; the latter guarantee amount supplies neither a liability
valuation nor a payment forecast. A dated parent debt-service path still needs
instrument terms, cash dates and permissible funding transfers
\citep[pp.~281--282, 310--311]{jpmorgan2025tenk}.

